\documentclass[a4paper,12pt]{article}
\usepackage[T2A]{fontenc}
\usepackage[utf8]{inputenc}
\usepackage[russian]{babel}
\usepackage{amsmath,amssymb,mathtools}
\usepackage{helvet}
\renewcommand{\familydefault}{\sfdefault}
\usepackage{icomma}
\usepackage[a4paper,left=19mm,right=19mm,top=19mm,bottom=21mm]{geometry}
\usepackage{microtype}
\usepackage{tikz}
\usetikzlibrary{calc}
\usepackage{tabularx,booktabs}
\usepackage{enumitem}
\usepackage{parskip}
\usepackage{fancyhdr}
\usepackage{setspace}
\usepackage{xcolor}
\definecolor{accent}{HTML}{008080}
\definecolor{darkaccent}{HTML}{004D4D}
\definecolor{textgray}{HTML}{333333}
\definecolor{softgray}{HTML}{707070}
\definecolor{linegray}{HTML}{D6DEDE}
\color{textgray}
\onehalfspacing
\setlength{\parindent}{0pt}
\setlength{\headheight}{15pt}
\setlength{\emergencystretch}{2em}
\setlist[itemize]{leftmargin=6mm,itemsep=2pt,topsep=4pt}
\setlist[enumerate]{leftmargin=8mm,itemsep=7pt,topsep=5pt}
\pagestyle{fancy}
\fancyhf{}
\lhead{\small\color{softgray} Степени и степенные преобразования}
\rhead{\small\color{softgray} Анастасия Павлова}
\cfoot{\small\color{softgray}\thepage}
\renewcommand{\headrulewidth}{0.4pt}
\newcommand{\checkitem}{\(\square\)\;}
\newcommand{\important}[1]{\textcolor{darkaccent}{\textbf{#1}}}
\begin{document}
\begin{titlepage}
\thispagestyle{empty}
\begin{center}
\vspace*{28mm}
{\Large\color{accent}\textbf{ЧЕК-ЛИСТ}}
\vspace{9mm}

{\fontsize{26}{31}\selectfont\bfseries Степени и степенные\\[2mm] преобразования выражений}
\vspace{8mm}

{\large Вычисления, корни, дробные и отрицательные показатели,\\ буквенные выражения}
\vspace{17mm}

{\color{darkaccent}\rule{55mm}{0.8pt}}
\vspace{17mm}

{\large\textbf{Подготовка к ЕГЭ}}
\vfill

{\large Лёвин Артём Александрович\\[2mm]\textbf{эксклюзивно для Анастасии Павловой}}
\vspace{8mm}

{\color{softgray}\today}
\vspace{23mm}
\end{center}
\begin{tikzpicture}[remember picture,overlay]
\draw[accent,line width=1.3pt]
($(current page.south west)+(18mm,15mm)$)
.. controls ($(current page.south west)+(63mm,26mm)$)
and ($(current page.south east)+(-68mm,6mm)$)
.. ($(current page.south east)+(-18mm,17mm)$);
\end{tikzpicture}
\end{titlepage}

\section*{1. Главная стратегия}
Почти любую вычислительную задачу со степенями удобно начинать с вопроса:
\begin{center}
\important{Можно ли привести составные числа, корни и дроби к степенному виду?}
\end{center}
\begin{itemize}
\item \checkitem составное число представить степенью удобного основания;
\item \checkitem корень заменить степенью с дробным показателем;
\item \checkitem десятичную дробь сначала перевести в обыкновенную;
\item \checkitem отрицательный показатель связать с обратным числом;
\item \checkitem степени одного основания собрать вместе;
\item \checkitem отдельно вычислить итоговый показатель.
\end{itemize}

\renewcommand{\arraystretch}{1.35}
\begin{table}[htbp]
\centering
\caption{Полезные переходы к степенному виду}
\begin{tabularx}{\linewidth}{@{}p{0.29\linewidth}X@{}}
\toprule
\textbf{Исходный объект} & \textbf{Полезная запись}\\
\midrule
\(49\) & \(49=7^2\)\\
\(8\) & \(8=2^3\)\\
\(18\) & \(18=2\cdot3^2\)\\
\(0{,}5\) & \(0{,}5=\dfrac12=2^{-1}\)\\
\(0{,}2\) & \(0{,}2=\dfrac15=5^{-1}\)\\
\(\sqrt[k]{a^m}\), \(a>0\) & \(\sqrt[k]{a^m}=a^{m/k}\)\\
\bottomrule
\end{tabularx}
\end{table}

\section*{2. Главные формулы}
\[
a^m\cdot a^n=a^{m+n},\qquad
\frac{a^m}{a^n}=a^{m-n}\quad(a\ne0).
\]
\[
(a^m)^n=a^{mn},\qquad (ab)^n=a^n b^n.
\]
\[
a^0=1\quad(a\ne0),\qquad a^{-n}=\frac1{a^n}\quad(a\ne0).
\]

\textbf{Пример:}
\[
3^{\sqrt5+10}\cdot3^{-5-\sqrt5}
=3^{(\sqrt5+10)+(-5-\sqrt5)}
=3^5=243.
\]

При делении сложный показатель знаменателя следует брать в скобки:
\[
\frac{a^p}{a^{q-r}}=a^{p-(q-r)}=a^{p-q+r}.
\]

\section*{3. Корни и дробные показатели}
Для положительного основания:
\[
\sqrt a=a^{1/2},\qquad \sqrt[3]{a}=a^{1/3},\qquad \sqrt[k]{a^m}=a^{m/k}.
\]

\textbf{Пример:}
\[
\frac{\sqrt[15]{5}\cdot\sqrt[10]{5}}{\sqrt[6]{5}}
=5^{\frac1{15}+\frac1{10}-\frac16}
=5^{\frac2{30}+\frac3{30}-\frac5{30}}
=5^0=1.
\]

Для корня чётной степени подкоренное выражение должно быть неотрицательным. Знаменатель не может обращаться в ноль. При едином степенном алгоритме с дробными показателями удобно отдельно фиксировать условие \(a>0\).

Например, при \(a>0\):
\[
\left(\frac{\sqrt a}{a^{1/3}}\right)^4
=\left(a^{1/2-1/3}\right)^4
=\left(a^{1/6}\right)^4
=a^{2/3}.
\]

\section*{4. Десятичные дроби и отрицательные показатели}
Десятичную дробь полезно убрать как можно раньше:
\[
0{,}5=\frac12=2^{-1},\qquad 0{,}2=\frac15=5^{-1}.
\]

Например,
\[
(0{,}5)^{\sqrt{10}+1}
=(2^{-1})^{\sqrt{10}+1}
=2^{-\sqrt{10}-1}.
\]

Для отрицательной степени:
\[
5^{-3}=\frac1{5^3}=\frac1{125}.
\]

\section*{5. Скобки и знаки}
Полезная памятка:
\begin{center}
\important{ВУПС: вычитаемое, умножаемое и подставляемое выражение --- в скобки.}
\end{center}
Особенно важно применять это правило, когда внутри выражения присутствуют знаки \(+\) или \(-\).

Например,
\[
p-(q-r)=p-q+r.
\]
Сначала раскрывайте скобки, затем приводите подобные слагаемые.

\section*{6. Буквенные выражения}
Если буквенная часть полностью совпадает, коэффициенты можно сложить:
\[
7m^{30}+11m^{30}=18m^{30}.
\]

При \(m\ne0\):
\[
\frac{7m^{30}+11m^{30}}{(3m^{15})^2}
=\frac{18m^{30}}{9m^{30}}=2.
\]

Сокращать можно одинаковые множители. Запись с двоеточием удобно заменять дробной чертой:
\[
18x^7\cdot x^{13}:(3x^{10})^2
=\frac{18x^7\cdot x^{13}}{(3x^{10})^2}.
\]

В выражении
\[
(3x^{10})^2=3^2(x^{10})^2=9x^{20}
\]
квадрат относится ко всему содержимому скобок.

\section*{7. Частые ошибки}
\begin{enumerate}[label=\textbf{\arabic*)}]
\item Формулу умножения степеней нельзя применять через знак сложения: \(a^m+a^n\ne a^{m+n}\) в общем случае.
\item Не теряйте внешнюю степень после упрощения содержимого скобок.
\item При возведении произведения в степень показатель относится ко всем множителям.
\item Проверяйте разложение чисел: \(18=2\cdot3^2\), а \(3^3=27\).
\item При вычитании сложного показателя сохраняйте скобки: \(p-(q-r)=p-q+r\).
\item Сначала упрощайте буквенное выражение, затем выполняйте числовую подстановку.
\item Разбивайте длинное преобразование на короткие действия для самопроверки.
\end{enumerate}

\section*{8. Алгоритм решения}
\begin{enumerate}[label=\checkitem]
\item Определены ограничения и ОДЗ.
\item Составные числа представлены удобными степенями.
\item Корни переведены в дробные показатели.
\item Десятичные дроби заменены дробями или степенями.
\item Сложные вычитаемые и подставляемые выражения заключены в скобки.
\item При умножении одинаковых оснований показатели сложены.
\item При делении одинаковых оснований показатели вычтены.
\item При возведении степени в степень показатели перемножены.
\item Степень произведения применена ко всем множителям.
\item Дробные показатели приведены к общему знаменателю.
\item Подобные слагаемые приведены.
\item Итоговый показатель перепроверен отдельно.
\item Подстановка выполнена после упрощения.
\item Ответ проверен с учётом исходных ограничений.
\end{enumerate}

\clearpage
\section*{Тренировочный блок}
\textbf{Путь Б: 04.09--06.09.}\par
На каждый день запланировано пять заданий. Для пятого задания включите таймер и запишите время выполнения.

\subsection*{04.09}
\begin{enumerate}[label=\textbf{\arabic*.}]
\item Найдите значение: \[3^{\sqrt5+10}\cdot3^{-5-\sqrt5}.\]
\item Найдите значение: \[\frac{(49^6)^3}{(7^7)^5}.\]
\item Найдите значение: \[\frac{5^{3\sqrt7-1}\cdot5^{1-\sqrt7}}{5^{2\sqrt7-1}}.\]
\item Найдите значение: \[\frac{8^{1-\sqrt7}}{2^{2-3\sqrt7}}.\]
\item Найдите значение: \[\frac{(0{,}5)^{\sqrt{10}+1}}{2^{-\sqrt{10}}}.\]
\end{enumerate}

\subsection*{05.09}
\begin{enumerate}[label=\textbf{\arabic*.}]
\item Найдите значение: \[\frac{\sqrt[15]{5}\cdot\sqrt[10]{5}}{\sqrt[6]{5}}.\]
\item При \(a>0\) упростите: \[\left(\frac{\sqrt a}{a^{1/3}}\right)^6.\]
\item При \(m\ne0\) упростите: \[\frac{7m^{30}+11m^{30}}{(3m^{15})^2}.\]
\item При \(x\ne0\) упростите: \[\frac{18x^7\cdot x^{13}}{(3x^{10})^2}.\]
\item При \(x\ne0\) найдите значение: \[\frac{(3x)^3\cdot x^{-9}}{2x^{-10}\cdot x^4}.\]
\end{enumerate}

\subsection*{06.09}
\begin{enumerate}[label=\textbf{\arabic*.}]
\item Найдите значение: \[(0{,}2)^2\cdot5^{-1}.\]
\item Упростите: \[\frac{\sqrt[12]{2^5}\cdot\sqrt[8]{2^3}}{\sqrt[6]{2}}.\]
\item При \(a\ne0\), \(b\ne0\) упростите: \[\frac{(4a^6b^{-2})^2}{8a^5b^{-1}}.\]
\item При \(x\ne0\) упростите: \[\frac{6x^{12}+9x^{12}}{(3x^6)^2}.\]
\item При \(a>0\) упростите \[\left(\frac{\sqrt a}{a^{1/3}}\right)^4,\] затем найдите значение при \(a=8\).
\end{enumerate}

\clearpage
\section*{Ответы к тренировочному блоку}
\renewcommand{\arraystretch}{1.35}
\begin{table}[htbp]
\centering
\caption{Ответы к заданиям 04.09--06.09}
\begin{tabularx}{\linewidth}{@{}p{0.19\linewidth}p{0.18\linewidth}X@{}}
\toprule
\textbf{Дата} & \textbf{Задание} & \textbf{Ответ}\\
\midrule
04.09 & 1 & \(243\)\\
04.09 & 2 & \(7\)\\
04.09 & 3 & \(5\)\\
04.09 & 4 & \(2\)\\
04.09 & 5 & \(\dfrac12\)\\
\midrule
05.09 & 1 & \(1\)\\
05.09 & 2 & \(a\)\\
05.09 & 3 & \(2\)\\
05.09 & 4 & \(2\)\\
05.09 & 5 & \(\dfrac{27}{2}=13{,}5\)\\
\midrule
06.09 & 1 & \(\dfrac1{125}\)\\
06.09 & 2 & \(2^{5/8}\)\\
06.09 & 3 & \(\dfrac{2a^7}{b^3}\)\\
06.09 & 4 & \(\dfrac53\)\\
06.09 & 5 & \(4\)\\
\bottomrule
\end{tabularx}
\end{table}
\end{document}
